\chapter*{Lampiran}
\addcontentsline{toc}{chapter}{Lampiran}

% ============================================================
% Lampiran A: Rubrik Penilaian
% ============================================================
\section*{Lampiran A: Rubrik Penilaian Tugas Praktik}

\begin{table}[h]
\centering
\footnotesize
\begin{tabular}{|l|l|l|l|}
\hline
\textbf{Kriteria} & \textbf{Sangat Baik} & \textbf{Baik} & \textbf{Perlu Perbaikan} \\
\hline
Desain Scene & Struktur jelas & Cukup jelas & Belum terorganisir \\
\hline
GameObject Setup & Komponen tepat & Sebagian tepat & Banyak kesalahan \\
\hline
Script Implementation & Logic benar & Ada kekurangan & Banyak kesalahan \\
\hline
Code Quality & Naming rapi & Ada duplikasi & Banyak duplikasi \\
\hline
Game Mechanics & Sesuai desain & Ada bug minor & Tidak berjalan baik \\
\hline
\end{tabular}
\caption{Rubrik Penilaian Tugas Praktik}
\end{table}

% ============================================================
% Lampiran B: Referensi Unity API
% ============================================================
\section*{Lampiran B: Referensi Unity API Penting}

\subsection*{GameObject Methods}
\begin{itemize}
  \item \texttt{GameObject.Find(string name)}: Mencari GameObject berdasarkan nama
  \item \texttt{GameObject.FindWithTag(string tag)}: Mencari GameObject berdasarkan tag
  \item \texttt{Instantiate(Object original)}: Membuat instance baru dari object
  \item \texttt{Destroy(Object obj)}: Menghapus object dari scene
\end{itemize}

\subsection*{Transform Methods}
\begin{itemize}
  \item \texttt{transform.position}: Posisi object dalam world space
  \item \texttt{transform.rotation}: Rotasi object
  \item \texttt{transform.localScale}: Skala object
  \item \texttt{transform.Translate(Vector3 translation)}: Menggerakkan object
  \item \texttt{transform.Rotate(Vector3 eulerAngles)}: Merotasi object
\end{itemize}

\subsection*{Component Methods}
\begin{itemize}
  \item \texttt{GetComponent\u003cT\u003e()}: Mendapatkan component dari GameObject
  \item \texttt{AddComponent\u003cT\u003e()}: Menambahkan component ke GameObject
  \item \texttt{GetComponentInChildren\u003cT\u003e()}: Mencari component di child objects
\end{itemize}

\subsection*{Input Methods}
\begin{itemize}
  \item \texttt{Input.GetKey(KeyCode key)}: Cek apakah key sedang ditekan
  \item \texttt{Input.GetKeyDown(KeyCode key)}: Cek apakah key baru ditekan
  \item \texttt{Input.GetMouseButton(int button)}: Cek apakah mouse button ditekan
  \item \texttt{Input.GetAxis(string axisName)}: Mendapatkan nilai input axis
\end{itemize}
